% ~250 words. Short scaffolding section. Compliance items; defers detail.
\section{Methods}
\label{sec:methods}

We describe FLAIR as three coupled artifacts: a data layer (CLIF, with
MIMIC-CLIF as the public development entry point), a task layer (five
benchmark definitions implemented as single-file Python modules), and
a protocol layer (containerized model-to-data execution with
aggregate-only JSON reporting). Detailed specification is deferred to
\S\ref{sec:datasets}--\S\ref{sec:protocols}; per-task algorithmic
schematics are in the Supplementary Appendix.

Reporting follows the TRIPOD+AI statement
\citep{collins2024tripodai}; the FLAIR library auto-populates the
checklist into the per-task \texttt{tripod\_ai.json} report.
A Datasheet for Datasets \citep{gebru2021datasheets} accompanies the
public MIMIC-CLIF cohorts, and a Model Card \citep{mitchell2019modelcards}
accompanies the baseline gradient-boosted-tree models.

\paragraph{Institutional review.} Each participating consortium site
obtained independent institutional review board (IRB) approval for
participation in CLIF and for execution of FLAIR benchmark code on its
patient data. No patient-level data are transmitted between sites; only
aggregate JSON outputs leave each institution. Site governance retains
final approval over which submitted models are executed.

\paragraph{Data availability.} MIMIC-CLIF is publicly accessible via
PhysioNet under the standard MIMIC data use agreement. Consortium data
are not redistributable; they are accessible only by submitting code
through the FLAIR portal for federated execution at participating sites.

\paragraph{Code availability.} The FLAIR library, task definitions,
report generators, and the CLIF-WorkBench container specification are
released under Apache~2.0 at \texttt{github.com/clif-consortium/FLAIR}.
